% 两页目标完成版：只在对应项目里程碑真实完成后转为投递版。
\logosection{\faUser}{个人简介}
东华大学软件工程 2027 届，聚焦 AI 应用与 Java 后端。主导教育 AI 学习平台并完成从手写 LLM/RAG 编排到 Spring AI 标准化架构的演进；另有已上线 AI 文字冒险游戏。具备 Prompt、RAG、Memory、Tool Calling、MCP、SSE 流式与全栈交付经验，能够从业务问题出发完成模型接入、状态管理、工具边界、观测评测和部署。

\logosection{\faCogs}{核心技能}
\begin{itemize}[parsep=0.35ex]
  \item \textbf{AI 应用架构}：Spring AI、ChatClient / ChatModel / Advisor、ChatMemory、RAG、Tool Calling、MCP；理解上下文工程、ReAct、模型无状态与 Agent 权限边界。
  \item \textbf{后端与数据}：Java 21、Spring Boot 4、MyBatis、MySQL、Redis、Qdrant；熟悉 REST、SSE、异常治理、指标观测、Docker、Nginx、Git。
  \item \textbf{Python / 前端}：Python、FastAPI、LlamaIndex；Vue 3、Next.js、TypeScript、微信小程序，可独立完成 AI 产品前后端闭环。
  \item \textbf{模型基础}：理解 Transformer / Attention；完成 BERT / LSTM 中文 NER 与 BART 文本生成微调实验。
\end{itemize}

\logosection{\faWrench}{项目经历}

\datedline{\textbf{408Master：面向 408 备考的 AI 学习平台}}{校实训二等奖 · 任队长}
\datedline{\tripleInfo{Java 21 / Spring Boot 4 / Spring AI 2}{Redis / Qdrant / LlamaIndex}{Vue 3 / 微信小程序}}{\href{https://github.com/tiankai-333/master408}{github.com/tiankai-333/master408}}

\textbf{项目背景：}完成题库、考试、错题、个人画像、知识库、AI 解析与组卷等产品闭环并上线后，继续治理手写 HTTP/SSE、正则路由、双轨 RAG 和一次性柏拉图教学等架构问题；以“保留成熟业务、替换 AI 技术边界”为原则持续演进。

\textbf{核心工作：}
\begin{itemize}[parsep=0.4ex]
  \item \textbf{渐进式迁移}：复用题库、组卷、Controller、DTO、MyBatis Mapper、Vue 与小程序接口，保留 legacy/Spring AI 双实现开关；以固定问题集对比功能正确性、p50/p95/p99、首 token、token 成本、错误率和改动面，避免重写业务或把框架升级误判为模型能力提升。
  \item \textbf{调用标准化}：以 ChatClient 取代手写请求和 SSE 解析，通过 ChatModel 适配 OpenAI 兼容模型；将 Provider 选择、Prompt、使用日志和模型传输拆分，Advisor 统一采集 token、费用、耗时和失败原因。
  \item \textbf{多轮教学}：使用 Redis ChatMemory 持久化会话，将一次性“柏拉图解析”改为逐轮追问；以用户与会话组合生成 conversationId，隔离不同学生上下文，并通过消息窗口与摘要控制上下文长度。
  \item \textbf{RAG 规范化}：由 LlamaIndex Python 边车负责 PDF/HTML 解析、语义分块、图片本地化与增量入库；Java 查询侧统一通过 Spring AI VectorStore 检索，保留引用展示及向量库异常时的无 RAG 降级。
  \item \textbf{真实工具调用}：复用已验证的错题查询和确定性组卷服务，将查题、错题优先组卷等能力封装为 Tool；模型只选择工具和参数，数据库查询、试卷生成及副作用仍由 Java 代码和权限规则控制。
  \item \textbf{受监督运维}：通过 MCP 暴露脱敏健康检查、日志、指标与 AI 使用统计等只读能力；清缓存、重建索引、重启等操作只生成待审批工单，不允许模型自动执行。
\end{itemize}

\textbf{关键工程决策：}
\begin{itemize}[parsep=0.35ex]
  \item 不以“用了 Spring AI”宣称 Agent：区分由系统预设执行的 Advisor 与由模型选择的 Tool，并为工具参数、调用次数、权限和终止条件设置边界。
  \item 不预设新框架性能更好：同时进行真实模型端到端测试和 Mock 隔离测试，区分网络/推理时间与框架开销；以可维护、可测试、可观测和后续扩展成本综合判断收益。
  \item API Key 从环境变量注入，数据库密文采用 AES-GCM；日志禁止输出明文或 key 前缀，公开仓库隔离旧 seed、jar、日志和构建产物。
\end{itemize}

\newpage
\noindent\textbf{吴天凯}\hfill AI 应用工程师 / Java 后端开发\\[-0.4ex]
{\color{ICBlue}\rule{\textwidth}{0.5pt}}

\logosection{\faWrench}{项目经历（续）}

\datedline{\textbf{StoryCraft：AI 文字冒险游戏（已上线）}}{}
\datedline{\tripleInfo{Python / FastAPI}{Next.js / TypeScript}{Agent Runtime / SSE}}{\href{https://github.com/baodashuaige/StoryCraft}{github.com/baodashuaige/StoryCraft}}

\textbf{项目背景：}构建可在线游玩的 AI 文字冒险游戏，核心问题不是单纯生成文本，而是让 NPC 对话自然、有记忆，同时不允许概率模型破坏确定性的世界规则和游戏状态。

\textbf{核心工作：}
\begin{itemize}[parsep=0.45ex]
  \item \textbf{Agent Runtime}：自建运行时维护场景、角色、线索、动作和结局状态，明确“LLM 负责叙事与候选动作、确定性代码负责校验和状态变更”的职责边界。
  \item \textbf{NPC 系统}：以人格、秘密、信任值、线索和剧情记忆构成角色上下文；不同玩家选择改变 NPC 反馈及后续分支，支持多结局叙事。
  \item \textbf{输出约束}：使用 Prompt 模板规定角色身份、可知信息和输出结构，并通过门控校验拦截越权动作、非法状态变化和破坏设定的内容。
  \item \textbf{全栈交付}：后端提供 REST API 与 SSE 流式输出，前端使用 Next.js / TypeScript；同时提供 Web 与 CLI 两种入口，API Key 在服务端加密保存。
  \item \textbf{质量与复用}：建立状态、动作、对话和接口单元测试，测试结果 176 passed；将 NPC 对话能力抽取为开源 npm 包 \texttt{@wutiankai/npc-dialogue}。
\end{itemize}

\textbf{项目结果：}产品已部署至 \href{https://game.mastertooling.shop}{game.mastertooling.shop}，具备可公网访问的完整游戏流程。

\datedline{\textbf{自然语言处理课程项目：NER 与文本生成}}{课程项目}
\datedline{\tripleInfo{BERT / LSTM}{BART}{CLUENER2020}}{}
\begin{itemize}[parsep=0.45ex]
  \item 基于 CLUENER2020 完成中文命名实体识别的数据预处理、标签编码、训练与评估，对比 BERT 与 LSTM 在上下文表征能力上的差异。
  \item 基于 BART 完成文本生成微调实验，理解编码器-解码器、Teacher Forcing、训练损失与生成阶段解码参数的作用。
  \item 通过课程实验补足 AI 应用工程之外的模型基础，能够区分模型训练、Prompt、RAG 与工具调用分别解决的问题。
\end{itemize}

\logosection{\faGithub}{开源与工程实践}
\begin{itemize}[parsep=0.45ex]
  \item 开源 npm 包 \texttt{@wutiankai/npc-dialogue}：从 StoryCraft 中抽取 NPC 对话能力，沉淀可复用接口与示例。
  \item 维护 408Master 的改造记录、概念笔记、阶段计划和前后对比文档；每个里程碑同时记录实现成本、性能/质量变化、收益与下一阶段决策。
  \item 使用 Git 管理改造基线与增量迁移，旧代码按“原样复用、机械升级、保留契约替换实现、淘汰”四档审计，减少无业务价值的重写。
\end{itemize}

\logosection{\faTrophy}{荣誉与证书}
\datedline{东华大学软件工程实训二等奖（408Master 项目队长）}{}
\datedline{人工智能训练师职业技能培训}{预计 2026.08 取证}

\logosection{\faGraduationCap}{教育经历}
\datedline{\textbf{东华大学（211）}}{\dateRange{2023.09}{2027.06}}
\datedline{\tripleInfo{软件工程}{本科 · 2027 届}{计算机科学与技术学院}}{上海}
\datedline{主修课程：数据结构、算法、操作系统、计算机网络、数据库、自然语言处理}{}
